\begin{lemma}
\label{lemma-quasi-coherence-higher-direct-images}
Let $f : X \to S$ be a morphism of schemes.
Assume that $f$ is quasi-separated and quasi-compact.
\begin{enumerate}
\item For any quasi-coherent $\mathcal{O}_X$-module $\mathcal{F}$ the
higher direct images $R^pf_*\mathcal{F}$ are quasi-coherent on $S$.
\item If $S$ is quasi-compact, there exists an integer $n = n(X, S, f)$
such that $R^pf_*\mathcal{F} = 0$ for all $p \geq n$ and any
quasi-coherent sheaf $\mathcal{F}$ on $X$.
\item In fact, if $S$ is quasi-compact we can find $n = n(X, S, f)$
such that for every
morphism of schemes $S' \to S$ we have $R^p(f')_*\mathcal{F}' = 0$
for $p \geq n$ and any quasi-coherent sheaf $\mathcal{F}'$
on $X'$. Here $f' : X' = S' \times_S X \to S'$ is the base change of $f$.
\end{enumerate}
\end{lemma}

\begin{proof}
We first prove (1). Note that under the hypotheses of the lemma the sheaf
$R^0f_*\mathcal{F} = f_*\mathcal{F}$ is quasi-coherent by
Schemes, Lemma \ref{schemes-lemma-push-forward-quasi-coherent}.
Using
Cohomology, Lemma \ref{cohomology-lemma-localize-higher-direct-images}
we see that forming higher direct images commutes with restriction
to open subschemes. Since being quasi-coherent is local on $S$ we
reduce to the case discussed in the next paragraph.

\medskip\noindent
Proof of (1) in case $S$ is affine. We will use the induction principle.
Since $f$ quasi-compact and quasi-separated we see that $X$
is quasi-compact and quasi-separated. For $U \subset X$
quasi-compact open and $a = f|_U$ we let $P(U)$ be the property that
$R^pa_*\mathcal{F}$ is quasi-coherent on $S$ for all quasi-coherent modules
$\mathcal{F}$ on $U$ and all $p \geq 0$. Since $P(X)$ is (1), it suffices
the prove conditions (1) and (2) of Lemma \ref{lemma-induction-principle} hold.
If $U$ is affine, then $P(U)$ holds because $R^pa_*\mathcal{F} = 0$ for
$p \geq 1$ (by  Lemma \ref{lemma-relative-affine-vanishing} and
Morphisms, Lemma \ref{morphisms-lemma-morphism-affines-affine})
and we've already observed the result holds for $p = 0$ in the first
paragraph. Next, let $U \subset X$ be a quasi-compact open, $V \subset X$
an affine open, and assume $P(U)$, $P(V)$, $P(U \cap V)$ hold.
Let $a = f|_U$, $b = f|_V$, $c = f|_{U \cap V}$, and $g = f|_{U \cup V}$.
Then for any quasi-coherent $\mathcal{O}_{U \cup V}$-module $\mathcal{F}$
we have the relative Mayer-Vietoris sequence
$$
0 \to
g_*\mathcal{F} \to
a_*(\mathcal{F}|_U) \oplus b_*(\mathcal{F}|_V) \to
c_*(\mathcal{F}|_{U \cap V}) \to
R^1g_*\mathcal{F} \to \ldots
$$
see Cohomology, Lemma \ref{cohomology-lemma-relative-mayer-vietoris}.
By $P(U)$, $P(V)$, $P(U \cap V)$ we see that
$R^pa_*(\mathcal{F}|_U)$, $R^pb_*(\mathcal{F}|_V)$ and
$R^pc_*(\mathcal{F}|_{U \cap V})$ are all quasi-coherent.
Using the results on quasi-coherent sheaves in
Schemes, Section \ref{schemes-section-quasi-coherent}
this implies that each of the sheaves
$R^pg_*\mathcal{F}$ is quasi-coherent since it sits in the middle of a short
exact sequence with a cokernel of a map between quasi-coherent sheaves
on the left and a kernel of a map between quasi-coherent sheaves on the right.
Whence $P(U \cup V)$ and the proof of (1) is complete.

\medskip\noindent
Next, we prove (3) and a fortiori (2). Choose a finite affine open
covering $S = \bigcup_{j = 1, \ldots m} S_j$. For each $j$ choose
a finite affine open covering
$f^{-1}(S_j) = \bigcup_{i = 1, \ldots t_j} U_{ji} $.
Let
$$
d_j = \max\nolimits_{I \subset \{1, \ldots, t_j\}}
\left(|I| + t(\bigcap\nolimits_{i \in I} U_{ji})\right)
$$
be the integer found in
Lemma \ref{lemma-vanishing-nr-affines-quasi-separated}.
We claim that $n(X, S, f) = \max d_j$ works.

\medskip\noindent
Namely, let $S' \to S$ be a morphism of schemes and let
$\mathcal{F}'$ be a quasi-coherent sheaf on $X' = S' \times_S X$.
We want to show that $R^pf'_*\mathcal{F}' = 0$ for $p \geq n(X, S, f)$.
Since this question is local on $S'$ we may assume that $S'$ is affine
and maps into $S_j$ for some $j$. Then $X' = S' \times_{S_j} f^{-1}(S_j)$
is covered by the open affines $S' \times_{S_j} U_{ji}$, $i = 1, \ldots t_j$
and the intersections
$$
\bigcap\nolimits_{i \in I} S' \times_{S_j} U_{ji} =
S' \times_{S_j} \bigcap\nolimits_{i \in I} U_{ji}
$$
are covered by the same number of affines as before the base change.
Applying
Lemma \ref{lemma-vanishing-nr-affines-quasi-separated}
we get $H^p(X', \mathcal{F}') = 0$. By the first part of the proof
we already know that each $R^qf'_*\mathcal{F}'$ is quasi-coherent
hence has vanishing higher cohomology groups on our affine scheme $S'$,
thus we see that $H^0(S', R^pf'_*\mathcal{F}') = H^p(X', \mathcal{F}') = 0$
by Cohomology, Lemma \ref{cohomology-lemma-apply-Leray}.
Since $R^pf'_*\mathcal{F}'$ is quasi-coherent
we conclude that $R^pf'_*\mathcal{F}' = 0$.
\end{proof}